\section{Background and the Capability Gap}
\label{sec:background}

\subsection{One Update, Several Representations}

Consider two replicas that start from the same file.  Replica A changes bytes
in one region and replica B changes another.  A conventional local filesystem
records two unrelated block updates.  A file-level synchronizer sees two
concurrent file versions and must either select one or ask a content-specific
merger.  Encryption and remote storage add further mappings between logical
content, local physical blocks, and opaque provider objects.

\sfs{} preserves causal structure below every surface and separates logical
fragment identity from its physical backend and cipher encoding.  The basic
state of a content stream is

\begin{equation}
  U = (\mathit{uuid}, e, \ell, VV, M), \qquad
  M[f] = d_f = (a_f,c_f),
\end{equation}

where $e$ is the fragment-size exponent, $\ell$ the logical length of the last
fragment, $VV$ a sparse version vector, and $M$ an ordered map from fragment
index $f$ to a causal dot containing replica alias $a_f$ and counter $c_f$.
This tuple is the replica-invariant logical head, not the complete physical
record.  A materialized record additionally contains selected locations and a
parent address used to walk local retained lineage.  Parent is omitted from
$U$ because replicas may store the same logical head at different addresses
and construct different physical parent chains during import, relocation, or
re-ciphering.  It is consequently excluded from logical record signatures and
transport identity.  Physical locations and cipher encodings are omitted for
the same reason: they may change without creating a new logical version.

This model separates three granularities that are often conflated.  A
byte-range is the API granularity of a read or write.  A fragment is the unit of
storage, transfer, and merge classification.  A unit is the scope of the
version vector and user-visible conflict.  Consequently, \sfs{} provides
byte-addressed I/O but fragment-granular causality.

\subsection{Design Requirements}

The representation is useful only if it satisfies the surrounding systems
requirements.  First, reconstructing the current unit must not replay its
history: the head map and selected fragments are materialized.  Second,
publication of a new map and catalog roots must be atomic.  Third, a sync pass
must transfer changed or missing fragments while retaining all concurrent
frontiers.  Fourth, the transport must not need content keys or plaintext.
Fifth, physical encoding and execution surface must be separable from logical
identity so that kernel, userspace, embedded, and browser-facing adapters can
share the format.

These requirements are not individually new.  \Cref{tab:positioning} therefore
positions representative systems by their architectural center rather than by
a checklist constructed around \sfs{}.

\begin{figure}[t]
  \centering
  \resizebox{\linewidth}{!}{%
  \begin{tikzpicture}[
    surface/.style={draw, rounded corners, align=center, minimum width=2.7cm,
                    minimum height=0.8cm, fill=RoyalBlue!8},
    core/.style={draw, rounded corners, align=center, minimum width=8.8cm,
                 minimum height=1.25cm, fill=ForestGreen!10, very thick},
    backend/.style={draw, rounded corners, align=center, minimum width=3.1cm,
                    minimum height=0.9cm, fill=black!4},
    link/.style={<->, >=Latex, thick}]
    \node[surface] (kernel) at (-4.5,2.1) {Linux kernel VFS};
    \node[surface] (fuse)   at (-1.5,2.1) {FUSE mount};
    \node[surface] (api)    at ( 1.5,2.1) {Rust core / C ABI};
    \node[surface] (wasm)   at ( 4.5,2.1) {experimental wasm VFS};

    \node[core] (state) at (0,0)
      {v12 fragment-causal state\\
       UUID + version vector + position$\rightarrow$dot map\\
       materialized head, retained history and strains, authenticated metadata};

    \node[backend] (raw)  at (-4.0,-2.0) {raw block device};
    \node[backend] (file) at ( 0.0,-2.0) {host file or memory image};
    \node[backend] (sync) at ( 4.0,-2.0) {one sync transport\\blind SaaS or live peer};

    \draw[link] (kernel) -- (state);
    \draw[link] (fuse) -- (state);
    \draw[link] (api) -- (state);
    \draw[link] (wasm) -- (state);
    \draw[link] (state) -- (raw);
    \draw[link] (state) -- (file);
    \draw[link] (state) -- (sync);
  \end{tikzpicture}%
  }
  \caption{One logical representation crosses execution boundaries above and
  physical or trust boundaries below.  The hosted service stores protocol
  objects rather than interpreting filesystem operations.}
  \label{fig:architecture}
\end{figure}

\Cref{fig:architecture} is the paper's organizing distinction.  The surfaces
do not create their own version models, and the hosted path does not introduce
a server-authoritative namespace.  They adapt one logical state to different
execution and persistence boundaries.

\begin{table}[t]
  \centering
  \small
  \setlength{\tabcolsep}{4pt}
  \caption{Architectural centers of representative neighboring systems.  The
    rows are not an ordering: several systems provide stronger guarantees than
    \sfs{} under their chosen threat model.}
  \label{tab:positioning}
  \begin{tabular}{@{}p{0.19\linewidth}p{0.22\linewidth}p{0.22\linewidth}p{0.26\linewidth}@{}}
    \toprule
    System & Native state/version unit & Replication and trust & Primary surface \\
    \midrule
    CVFS~\citep{Soules2003CVFS} & per-write file versions & local storage & mounted versioning FS \\
    Ori~\citep{Mashtizadeh2013Ori} & immutable objects and commits & optimistic peer replication & mounted replicated FS \\
    Coda~\citep{Kistler1992Coda} & file and volume state & disconnected clients, trusted servers & mounted distributed FS \\
    \shortstack[l]{SUNDR / SPORC\\\citep{Li2004SUNDR,Feldman2010SPORC}} & signed version structures / operations & untrusted server with fork guarantees & repository / collaborative app \\
    WNFS~\citep{WNFSspec} & encrypted content-addressed nodes & state-based CRDT over untrusted storage & library and web applications \\
    Peergos~\citep{Peergos} & encrypted chunked filesystem & encrypted cloud/P2P synchronization & web and desktop clients \\
    Ceph~\citep{Weil2006Ceph} & distributed objects & server-coordinated trusted cluster & file, block, and object services \\
    \midrule
    \textbf{SFS} & position-addressed fragment version map & optimistic VV sync through HbC store or live peer & kernel/FUSE FS, library, experimental wasm \\
    \bottomrule
  \end{tabular}
\end{table}

WNFS is the closest web-facing neighbor: it makes encrypted,
content-addressed tree state a state-based CRDT for library and browser use,
whereas \sfs{} makes an ordered position-to-dot map the unit of local lineage
and reconciliation and carries that representation through mounted raw-device
and FUSE surfaces as well as library and experimental wasm surfaces.  This is
a difference in architectural center, not a claim that either design subsumes
the other.

The broader distinction is the role assigned to the fragment map.  CVFS and
snapshot filesystems establish that native versioning is useful.  Optimistic filesystems
establish version vectors and disconnected operation.  WNFS and Peergos show
that encrypted collaboration can be portable.  Ceph shows that a substrate
can serve several APIs.  \sfs{} combines lessons from these lines, but studies
the fragment map itself as the shared boundary between history,
synchronization, conflict isolation, and execution backends.  The NoSQL annex
tests semantic projection as a consequence, not as the organizing goal.
